\section{Results}\label{sec:results}

Of \TestSpecified{} hypotheses, \TestSupported{} are supported, \TestContradicted{} are
rejected with the effect running against the direction stated, and \TestNotRejected{} are not
distinguishable from no difference. Table~\ref{tab:descriptives} gives the descriptive
picture and Table~\ref{tab:tests} the tests. The sections below take them in the order that
makes the pattern visible, rather than in hypothesis order.

\input{generated/tables/descriptives}

\subsection{Liquidity supply falls and execution gets more expensive}

The clearest result is that the window is harder to trade in. The quoted spread widens from
\SpreadMeanBpsControl{} to \SpreadMeanBpsExpiry{} basis points, displayed notional in the
book falls, and a sweep of the displayed book by a fixed order moves the price further:
\MoveBpsSmallControl{} basis points per \ProbeSizeCr{} crore on a control session against
\MoveBpsSmallExpiry{} on an expiry. Per-trade price impact rises with it, and so does the
impact of a given quantity of signed order flow, which is the largest effect among the
liquidity measures ($d = \PFourEffectSizeCohenD{}$).

Figure~\ref{fig:cost} shows where this concentrates. The cost of moving the price is several
times higher in the illiquid group on any day, which is what one would expect and is worth
stating quantitatively rather than assuming; the expiry-day increase appears in both groups.

\begin{figure}[H]
  \centering
  \includegraphics[width=0.92\textwidth]{marking_cost.pdf}
  \caption{Price move, in basis points, from sweeping the displayed book by
    \ProbeSizeCr{}~crore, by group and session type. The sweep assumes the book neither
    replenishes nor reacts, so the figure is a lower bound on what moving the price would
    cost.}
  \label{fig:cost}
\end{figure}

\subsection{But prices are not more directional}

Every measure built to detect sustained pressure in one direction fails to find it, and the
point estimates are near zero or the wrong sign.

The variance ratio \emph{falls} from \VarianceRatioControl{} to \VarianceRatioExpiry{}
($d = \POneEffectSizeCohenD{}$, $p = \POnePValue{}$). Prices on expiry sessions are closer to
a random walk than on control sessions, not further from one. The autocorrelation of signed
order flow is unchanged ($d = \PTwoEffectSizeCohenD{}$, $p = \PTwoPValue{}$), so flow is no
more persistent. The permanent share of price impact is unchanged ($p = \PThreePValue{}$),
so the larger moves that flow produces are no more likely to survive. Submitted order flow is
no more one-sided ($p = \HOneSixPValue{}$).

Figure~\ref{fig:pressure} shows the distributions. These are the measures that separate the
two explanations, and they point the same way.

\begin{figure}[H]
  \centering
  \includegraphics[width=\textwidth]{pressure_measures.pdf}
  \caption{The three measures that distinguish directional pressure from heavier two-sided
    trading, by session type. Under sustained pressure the variance ratio rises above one,
    flow autocorrelation rises, and the permanent share of impact rises. None of them does.}
  \label{fig:pressure}
\end{figure}

\subsection{The settlement price ends up further from the close}

The settlement VWAP sits \AbsReversalBpsExpiry{} basis points from the closing price on
expiry sessions against \AbsReversalBpsControl{} on controls, a difference of
$d = \STwoEffectSizeCohenD{}$ and one of the stronger results in the study. The final
minute's VWAP also diverges further from the window's ($d = \SThreeEffectSizeCohenD{}$). The
gap against the mid at the window's open is larger too, but not reliably so
($p = \SOnePValue{}$).

This is the result that most resembles a settlement being pushed and then released. It is
also what a thinner, more volatile book produces on its own: realised variance per minute is
higher on expiry, and a noisier price series around an average will land further from any
single point of it. The two are not separable from these three numbers alone, which is why
the direction measures above matter---and they do not support the first reading.

Two findings cut against the pushing story more directly. Less volume transacts within a
basis point of the running VWAP on expiry sessions ($d = \SFourEffectSizeCohenD{}$), which is
the opposite of what benchmark-driven accumulation would produce. And volume does \emph{not}
concentrate into the final minute; it spreads out, against the stated direction
($d = \SFiveEffectSizeCohenD{}$, $p = \SFivePValue{}$). If the cheapest moment to move an
average is its end, the flow is not going there.

\begin{figure}[H]
  \centering
  \includegraphics[width=\textwidth]{settlement_benchmarks.pdf}
  \caption{Distance of the settlement VWAP from the mid at the window's open, from the
    closing price, and from the final minute's VWAP, by session type. Absolute values: the
    sign depends on which way a given participant is positioned and averages out across
    securities.}
  \label{fig:settlement}
\end{figure}

\subsection{Size is concealed differently, not more}

The share of resting size that is never displayed is high on any day---\ConcealedDepthShareControl{}
per cent of resting quantity on control sessions---and it does not move reliably on expiry
($p = \COnePValue{}$). What changes is the shape of concealment rather than its level.
Concealed parent orders are larger relative to displayed ones
($d = \CTwoEffectSizeCohenD{}$), and more entered orders carry a disclosed quantity below
their total ($d = \HNinePValue{}$ for the count measure). Meanwhile the replenishment rate
per unit matched \emph{falls}, against the stated direction
($d = \CThreeEffectSizeCohenD{}$).

Fewer, larger concealed parcels that are refilled less often per share traded is a
description of bigger institutional orders working through the window, not of size being
sliced to hide an accumulation.

\begin{figure}[H]
  \centering
  \includegraphics[width=0.92\textwidth]{concealment.pdf}
  \caption{Concealed size, by session type. The level barely moves; the size of individual
    concealed orders relative to displayed ones does.}
  \label{fig:concealment}
\end{figure}

\subsection{Order handling: more of everything, not more churn}

Cancellation counts rise sharply---the largest effect in the study
($d = \HEightEffectSizeCohenD{}$)---and so does the share of orders cancelled within a second
of entry. But cancellations \emph{per entry} fall, against the stated direction
($d = \HSevenEffectSizeCohenD{}$). Entries must therefore be rising faster than
cancellations: the window attracts more order submission of every kind rather than a
disproportionate amount of placing and pulling. Book pressure is also less persistent, not
more ($d = \HOneEightEffectSizeCohenD{}$), which is what happens when more participants are
active on both sides.

Proprietary desks take a larger share of traded volume, while algorithmic flow's share does
not change reliably ($p = \HFivePValue{}$).

\subsection{Within the window}

Figure~\ref{fig:profile} plots the main measures minute by minute. The expiry and control
lines are separated for most of the window rather than diverging at its end, which is the
shape one expects from a market that is generally busier and thinner rather than from
activity aimed at the closing average.

\begin{figure}[H]
  \centering
  \includegraphics[width=\textwidth]{window_profile.pdf}
  \caption{Spread, traded volume share, concealed share of resting size and cancellations per
    entry by minute from \DesignWindowStart{}, split by session type and group. Each line
    averages over securities and months.}
  \label{fig:profile}
\end{figure}

\subsection{Effect sizes}

Figure~\ref{fig:forest} places all the tested effects on one axis with their intervals. With
\TestPairsTotal{} matched pairs the tests have power to detect differences well below
anything economically interesting, so the ordering of effect sizes is more informative than
the pattern of $p$-values.

\begin{figure}[H]
  \centering
  \includegraphics[width=\textwidth]{effects_forest.pdf}
  \caption{Paired effect sizes with 95\% intervals. Filled markers are supported; open
    markers are rejected with the effect running against the stated direction.}
  \label{fig:forest}
\end{figure}

\input{generated/tables/tests}

\subsection{Is it the settlement, or is it the Thursday?}

Table~\ref{tab:contrasts} reports the two difference-in-differences described in
Section~\ref{sec:method}. \ContrastPlaceboNote{}

\input{generated/tables/contrasts}
